% NOTE (migration 23 Nov 2026): Text moved unchanged from main.tex. Much of it was drafted with AI assistance; rewrite it in your own words before submission (the supervisor runs AI checks).
% NOTE (migration 23 Nov 2026): This chapter still describes a PROPOSED study in the future tense, with [TBD] items. Rewrite it as the method actually followed, once the research-question decisions in RQ-EVIDENCE.md are agreed; mention the fine-tuning method (now in Implementation) and remove the statement that fine-tuning is out of scope.
\hypertarget{methodology-proposed-study}{%
\chapter{Methodology}\label{methodology-proposed-study}}

\begin{quote}
\textbf{Updated 25 October 2026 for supervisor review. This chapter describes a
PROPOSED research study. No AI extraction experiments have been run, no annotated research dataset has been completed, and no participants have been recruited. Synthetic scheduling fixtures and local engineering verification exist; they are reported separately from research results.}
Present-tense descriptions of procedures describe the plan.
Square-bracket items (\texttt{{[}TBD\ …{]}}) are decisions or
measurements not yet available. Options are proposed with justification
and need the owner\textquotesingle s and supervisor\textquotesingle s
confirmation (SRS Section\textasciitilde18). The chapter is aligned with
the revised SRS (DRAFT v0.2, 10 October 2026; 64 functional requirements); every evaluated capability maps to a specified
feature or an explicit research procedure (Section\textasciitilde3.19).
\end{quote}

\hypertarget{research-design-and-justification}{%
\section{Research design and
justification}\label{research-design-and-justification}}

\textbf{Design.} A design-and-evaluate study with three components:

\begin{enumerate}
\def\labelenumi{\arabic{enumi}.}
\tightlist
\item
  \textbf{An offline extraction experiment} (RQ1, RQ2, RQ4): candidate
  LLMs are run on a fixed, annotated set of requirements under
  equivalent conditions, and their raw outputs are scored against gold
  annotations. This is a controlled benchmark; it isolates extraction
  quality from human behaviour and from the solver.
\item
  \textbf{A human-in-the-loop comparison} (RQ3): manual form entry and
  AI-assisted entry with review are compared on equivalent tasks using
  the same scheduling engine. Because participant numbers are uncertain
  (OQ-12), it is planned as a \emph{pilot}, and its conclusions are
  limited accordingly.
\item
  \textbf{Engineering verification of the surrounding system}: tests
  that approval gates, the independent validator, solver-outcome
  reporting and publication gates behave as specified, and measurements
  of scheduling behaviour on synthetic instances. These use
  deterministic test oracles, not statistical inference.
\end{enumerate}

\textbf{Why this design.} The approved contribution is the extraction
method and its evaluation (project specification,
Section\textasciitilde3). Offline benchmarking is the established way to
compare models on a labelled task and permits repeated runs, paired
statistics and cost accounting. A human study is needed because the
practical claim is \emph{reduced effort without loss of correctness},
which a benchmark cannot show. Engineering verification is needed
because the prototype\textquotesingle s value depends on the claim that
unreviewed or invalid content cannot reach the solver or publication,
which is a property of the system, not of the model. Keeping the three
separate follows the principle that valid output is not correct output,
and that a valid schedule is not evidence of correct interpretation.

\textbf{What the study does not do.} It does not claim to evaluate
solver quality against other timetabling systems or to show that a
smaller prototype equals UniTime\textquotesingle s coverage. It does not
attribute differences in schedule quality to the AI when they arise from
different formal rules or objective weights (project specification,
Section\textasciitilde13).

\begin{longtable}[]{@{}p{\dimexpr(\linewidth-6\tabcolsep)/3\relax}p{\dimexpr(\linewidth-6\tabcolsep)/3\relax}p{\dimexpr(\linewidth-6\tabcolsep)/3\relax}@{}}
\toprule\noalign{}
RQ & Component & Primary evidence \\
\midrule\noalign{}
\endhead
\bottomrule\noalign{}
\endlastfoot
RQ1 extraction accuracy & Offline experiment & Raw and post-review
metrics on held-out data (Section\textasciitilde3.9) \\
RQ2 clarification / unsupported & Offline experiment & Clarification and
unsupported-request metrics (Section\textasciitilde3.9.3) \\
RQ3 effort vs manual & Human pilot & Time, corrections, final accuracy
(Section\textasciitilde3.10) \\
RQ4 model comparison & Offline experiment & Accuracy, latency, cost,
variability (Section\textasciitilde3.7, Section\textasciitilde3.9,
Section\textasciitilde3.12) \\
Supporting engineering question & Verification & Test outcomes
(Section\textasciitilde3.11) \\
\end{longtable}

\section[Users, needs and workflow]{Users, educational needs and proposed functional workflow}\label{system-architecture-and-separation-of-concerns}
The primary user is the timetable administrator. Instructors supply availability and preferences through the administrator in the initial scope; instructors and students view the published timetable within their authorised scope. Direct instructor submission is an optional extension. A researcher evaluates interpretation and task effort using approved study records. These roles should not be assumed to have identical access.

Before constructing the evaluation tasks, the supervisor and, where feasible, staff with timetabling experience will review the proposed needs: access to required activities, avoidance of group and staff clashes, suitable teaching facilities, feasible availability, and understandable reporting of incomplete or conflicting schedules. This consultation is planned, not completed. Any agreed requirement must fit the supported scope; new capabilities require an explicit scope decision.

The prototype is intended to support the following workflow:
\begin{enumerate}
\item The administrator records courses, teaching activities, instructors, student groups, rooms and teaching periods, including eligible instructors per activity component, mandatory weekly instructor workload limits, and strict or preferred daily limits for instructors and groups.
\item The administrator enters a scheduling need, either through a manual form or in ordinary language. The system displays its proposed interpretation and identifies missing or ambiguous information.
\item The administrator corrects or clarifies the proposal and approves it before it affects a timetable. For example, ``Dr Ahmed cannot teach on Tuesdays and prefers mornings'' represents a mandatory restriction and a separate preference; the institution must define ``morning''.
\item The system assigns eligible instructors as well as times and rooms, retaining the same instructor across weekly occurrences of each activity, and reports which activities were scheduled, whether mandatory checks passed, and which preferences remain unmet.
\item The administrator compares alternatives or changes a placement; instructor, time and room assignments may be locked. Each changed revision is checked again before publication. Changes in approved inputs require validation against current requirements.
\item A timetable is published only when the exact revision is complete and independently validated against current approved inputs, with every supported mandatory rule evaluated and passed. Missing data needed for a mandatory check blocks publication; unevaluated preferences are disclosed. Mandatory checks cannot be bypassed. Authorised instructors and students can then view the published schedule.
\end{enumerate}
This is the target workflow; its implemented subset and remaining work are recorded in Section~\ref{implementation-status}. The detailed acceptance criteria remain in the SRS. Generation, interpretation and review are evaluated separately so that a valid timetable cannot conceal a misread administrative requirement.

\subsection{Exceptions and recovery in the administrative workflow}
The revised SRS adds six supporting requirements (FR-059--FR-064) to protect the workflow when operations fail or users act on outdated information. These describe the complete target behaviours. Safeguards for implemented operations have local test evidence, but coverage of future extraction, revision and publication workflows remains incomplete (Section~\ref{implementation-status}).

\begin{itemize}
\item \textbf{Conflicting changes (FR-059).} If another administrator has changed a record, an outdated save is refused and the user is asked to reload and review it. Publication rechecks the revision, current inputs, validation and permissions at the point of commitment. A timetable generated before an input change may be published only after it passes validation against current inputs and satisfies every other publication condition; its original generation version remains visible.
\item \textbf{Interrupted operations (FR-060).} A lost connection must preserve acknowledged saves and allow the administrator to retrieve the operation status before trying again. An interrupted worker job is recorded as failed with a reason. An explicit retry creates a distinct job from the original snapshot unless the administrator chooses new inputs. An interrupted import must leave no partially committed dataset.
\item \textbf{Repeated requests and provider failures (FR-061).} Repeated delivery of the same request must not create duplicate saves, approvals or publications. An explicit new run remains a separate operation. Provider retries are bounded and recorded; failure preserves the statement and offers an explicit retry or manual entry. A late AI response cannot replace approved constraints.
\item \textbf{Incomplete results and cancellation (FR-062).} Where partial diagnostic results are available, the system identifies unscheduled activities and prevents publication of an incomplete timetable. Requesting cancellation and confirming it are separate states; simultaneous completion and cancellation must produce one accurate final outcome. Any retained partial export is labelled as a draft; the source of partial diagnostic timetables remains undecided. After cancellation, an administrator may explicitly retain a complete result only if it passes validation; it must disclose that the search stopped without proven optimality.
\item \textbf{Failed validation (FR-063).} A validator error, timeout or missing required check cannot count as a pass. The administrator receives a reason and a retry route, with prior reports retained. A later failed or incomplete check of the same revision and inputs blocks reliance on an earlier passing report. Unchecked preferences are not reported as having zero penalty.
\item \textbf{Expired sessions and changed permissions (FR-064).} A protected operation is refused when the session expires or access is revoked. After sign-in, current permissions and versions are checked again; approval or publication is not automatically replayed. Restricted timetable and study data remain inaccessible to unauthorised users.
\end{itemize}
These safeguards support reliable administrative decisions. They do not extend the educational scope to examinations, individual enrolments or other optional features.

\section{Supported requirements and interpretation criteria}\label{formal-representation-of-supported-constraints}
The initial requirement set covers instructor, room and group unavailability; permitted activity times; fixed placements; and room capacity, type and equipment. Proposed preferences include teaching periods, avoiding late sessions, reducing gaps and room changes, and balancing daily load. Daily limits for instructors and student groups are approved scope, with mandatory or preferred classification. Instructor eligibility, continuity across weekly occurrences and locked instructor assignments are required. Configured weekly instructor load caps are mandatory and measured in elapsed teaching hours, excluding breaks. A missing cap is disclosed as unconfigured rather than checked; zero prohibits assignments. The implemented daily policy uses hard global hour caps, activity-day distribution, consecutive-period caps and internal idle-period caps. The demonstration caps are six teaching hours per day for both students and instructors; these values are owner-selected settings, not an established institutional policy. Preferred or entity-specific daily-limit variants remain unfinished. Structural checks, including complete assignment and avoidance of room, instructor and group overlaps, remain active irrespective of how a requirement was entered.

The educational meaning of each requirement will be documented before annotation: whose activities are affected, whether it is mandatory or preferred, and which period, threshold or facility is intended. Terms such as ``late'', ``balanced'' and ``a long gap'' require agreed definitions. A proposal may remain incomplete while clarification is pending; approval requires the intended person or activity and all necessary details to be resolved. Technical field definitions are maintained in the project's constraint contract rather than treated as the main contribution of this chapter.

For evaluation, an interpretation is correct when it preserves the intended requirement type, affected entities, mandatory or preferred status, and parameter values. Equivalent wording and ordering do not count as errors. Preference weights are compared only when explicitly stated; otherwise the agreed profile defaults apply. This scoring rule and the semantic definitions will be fixed before running the evaluated models. The following sections retain the experimental details needed to make those measurements reproducible.

\hypertarget{dataset-construction}{%
\section{Dataset construction}\label{dataset-construction}}

Statement families will cover all four language categories, with a category label on each item. Translations and cross-language paraphrases expressing the same constraints belong to one family and must remain in the same split. An initial unreviewed synthetic release now contains 128 statements in 32 translated intent families: 96 development statements and 32 reserved statements. This is draft material rather than a completed, adjudicated research dataset; final study sizes remain undecided and no extraction model has been evaluated. The annotation guide must cover Arabic and Egyptian Arabic negation and hedging, mixed-script entity mentions and transliterated names. Where available, a second annotator fluent in Egyptian Arabic should review the relevant subset; availability is not assumed.

\textbf{Current draft release (25 October 2026).} Version \texttt{0.1.0-draft} contains 32 statements in each of the four language categories. Each family has four translations; 24 families are development material and eight are reserved. One fictional department context supplies instructors and aliases, rooms, groups, activities, the five named teaching periods and a fictional validated source draft for lock interpretation. The statements and proposed labels were authored with project-assistant support; they have not been independently annotated or adjudicated. They include clear and compound requests, negation, missing values, ambiguous and unknown entities, hard contradictions, unsupported requests, temporal boundaries and session-to-hour conversion.

The reserved material is \emph{not a blind test set}: its author has seen the labels and its source remains accessible. Default tooling does not read it, and explicit access requires a logged reason, but this is a local workflow guard rather than operating-system isolation. A fresh independently authored and adjudicated confirmatory set is required if this draft informs prompt or implementation choices. The proposed confirmatory protocol below has not been executed.

\textbf{Implemented offline scoring contract.} The contract distinguishes resolved proposals, clarification, unsupported content, conflicts and mixed requests. It covers typed manual rules, weekly instructor-hour updates, global daily-cap patches and assignment-lock proposals. Proposals never confer automatic approval. Exact whole-case scoring compares the decision, complete proposal multiset, clarification codes and unsupported codes; arbitrary rule titles and proposal ordering are ignored, while duplicate proposals count as extra predictions. Missing outputs fail the case, invalid attempted proposals count as false positives, and empty metric classes remain undefined rather than being credited as perfect. A conditional field-bag diagnostic does not replace strict whole-proposal or whole-case correctness.

The scorer also reports language breakdowns, cross-language family consistency, unknown entity references and available latency/token/attempt/cost metadata. Consistency can be high for consistently wrong outputs. Missing telemetry remains unknown. Raw and operator-declared post-review files are scored separately; the tool does not verify that human review occurred or measure correction effort. Gold-free exports contain only input text, identifiers and a context companion. Manifest, configuration, prediction and scorer hashes support reproducible rescoring. Oracle fixtures that copy expected labels are software self-tests only, not model-performance observations.

\textbf{Unit and content.} The unit is a \emph{requirement statement}
(one or more sentences) paired with an \emph{institutional context}
(entity tables: instructors with aliases, courses/activities, rooms,
groups, named periods) and a \emph{gold annotation}: the set of gold
constraints, the expected open items (clarification needed, unsupported
elements), and a category label.

\textbf{Categories} (a statement can carry several tags): simple (single
rule); compound (several rules); negated or polarity-sensitive
("cannot", "prefers no", "would rather not"); numeric (times and counts
in varied formats); ambiguous (entity or parameter underspecified);
undefined-term ("morning", "late" without a definition); contradictory
(within the statement or against an existing rule); unsupported (outside
the vocabulary); and paraphrase variants of the above.

\textbf{Sources of statements (to be agreed, OQ-09).} Options, which can
be combined: (a) statements authored by the researcher from the
constraint examples in the supplied literature (for example the exam and
constraint lists in~\cite{abdennadher_guc_exams}) and the candidate vocabulary; (b) statements elicited
from scheduling staff or instructors, if access and approval are
available; (c) wording drawn from real department communications,
anonymised, if permitted; (d) model-assisted paraphrasing of
human-authored statements, always reviewed by a human --- LLM-generated
statements are \emph{not} used as gold-labelled test material without
review because they could favour the same model family. Provenance of
every item is recorded.

\textbf{Institutional contexts.} The confirmatory study proposes at least two synthetic department
contexts with different names, aliases and period definitions
so that results do not depend on one entity table. {[}TBD: number and
size.{]} Entity names are fictional or anonymised.

\textbf{Size --- options, not decisions.} The project specification
proposes 150--300 requirements, subject to annotation effort. The table
shows how the 95 \% Wilson confidence-interval half-width for a
proportion shrinks with the number of \emph{independent} test statements
(illustrative arithmetic; not a target):

\begin{longtable}[]{@{}p{\dimexpr(\linewidth-8\tabcolsep)/4\relax}p{\dimexpr(\linewidth-8\tabcolsep)/4\relax}p{\dimexpr(\linewidth-8\tabcolsep)/4\relax}p{\dimexpr(\linewidth-8\tabcolsep)/4\relax}@{}}
\toprule\noalign{}
Test statements (n) & Observed accuracy 0.80 & Observed accuracy 0.90 &
Observed accuracy 0.95 \\
\midrule\noalign{}
\endhead
\bottomrule\noalign{}
\endlastfoot
100 & ± 7.8 points & ± 6.0 & ± 4.5 \\
150 & ± 6.4 & ± 4.8 & ± 3.7 \\
200 & ± 5.5 & ± 4.2 & ± 3.1 \\
\end{longtable}

Per-category estimates are wider because each category has fewer items;
claims about individual categories will be reported with their own
intervals and treated as exploratory unless enough items exist. Two
options for confirmation: \textbf{Option A} --- about 150 statements
(for example roughly 50 development, 100 held-out); \textbf{Option B}
--- about 300 statements (for example roughly 100 development, 200
held-out). A small \textbf{pilot set} of {[}TBD{]} statements, separate
from both, is used to debug the annotation guide and the pipeline. Final
numbers depend on annotation effort and on the pilot. Paraphrases of one
base statement are not independent; the effective n is the number of
statement \emph{families} (Section\textasciitilde3.6), and analysis will
be clustered accordingly.

\hypertarget{annotation-protocol-disagreement-and-agreement}{%
\section{Annotation protocol, disagreement and
agreement}\label{annotation-protocol-disagreement-and-agreement}}

\begin{enumerate}
\def\labelenumi{\arabic{enumi}.}
\tightlist
\item
  \textbf{Annotation guide} (written first, versioned): constraint
  definitions with examples and counter-examples; rules for resolving
  vague terms; what counts as unsupported; how to annotate compound,
  negated and contradictory statements; the expected clarification for
  each ambiguity type; canonicalisation rules
  (Section\textasciitilde3.3). The guide is developed on the pilot set
  and is versioned. The version used to annotate the held-out set is
  frozen before held-out annotation begins and is recorded in the
  held-out manifest. Development annotations may later be revised under
  newer guide versions; every revision is logged with its guide version
  and reason (Section\textasciitilde3.6).
\item
  \textbf{Annotators.} The researcher will independently annotate or review all research items; the current assistant-authored draft labels do not fulfil this obligation. A
  \textbf{second annotator} (to be identified; a supervisor-nominated
  person or fellow student with scheduling knowledge) independently
  annotates a subset {[}TBD size, preferably the full held-out set if
  feasible{]} blind to the first annotation and to any model output.
\item
  \textbf{Agreement measurement.} For categorical decisions (constraint
  type, severity, per-item "clarification needed", "unsupported")
  Cohen\textquotesingle s κ with confidence intervals; for entity and
  parameter fields, exact-match agreement; for whole statements, the
  proportion with identical canonical constraint sets. Because
  categories may be unbalanced, κ is reported with raw percentage
  agreement. Targets for acceptable agreement are not set in advance.
  {[}TBD with supervisor.{]}
\item
  \textbf{Adjudication.} Disagreements are discussed, the guide is
  clarified where the disagreement shows an ambiguity, and a final gold
  label is recorded with the rule used. A guide revision that changes
  development labels creates a new development-set version, and earlier
  development results remain attributed to the version they used.
  Held-out gold labels are changed only for clear annotation errors,
  decided without looking at any model output, logged, and applied to
  all models identically before the confirmatory run; any change after
  that run is reported as a deviation.
\item
  \textbf{Single-annotator limitation.} If a second annotator is
  unavailable, results are reported as relative to a single
  annotator\textquotesingle s interpretation and the limitation is
  stated prominently (Section\textasciitilde3.16).
\end{enumerate}

\hypertarget{development--validation--held-out-separation-and-leakage-prevention}{%
\section[Data separation and leakage]{Development / validation / held-out separation and leakage
prevention}\label{development--validation--held-out-separation-and-leakage-prevention}}

\begin{itemize}
\tightlist
\item
  \textbf{Splits.} A pilot set (tooling and guide refinement), a
  development set (prompt, example and schema-description iteration), an
  optional validation set (choice among frozen candidate configurations),
  and a held-out test set (confirmatory evaluation only).
\item
  \textbf{Held-out set: proposed frozen and isolated protocol, not yet executed.} Language-category representation will be determined before sealing, and leakage checks will compare translations across categories as well as within-language paraphrases. Held-out statement families
  are authored, annotated under the recorded guide version and adjudicated
  before any prompt development begins. They are then sealed: a manifest
  of item identifiers and content hashes is recorded, and the items are
  stored separately from development material. The experiment harness
  cannot load held-out items in development runs, refuses held-out runs
  unless every element of the configuration is frozen, and logs every
  held-out access (FR-049). After sealing, families are never added to,
  removed from or moved into the held-out set; corrections follow
  Section\textasciitilde3.5, item 4. Because the researcher authored the
  held-out items, this isolation limits, but cannot remove, the
  researcher's knowledge of them (Section\textasciitilde3.15).
\item
  \textbf{Development material: allowed to evolve, with versioned
  records.} Development statements, few-shot examples, prompt templates,
  schema descriptions and decoding settings may be added, revised or
  removed throughout development. Each change produces a new version of
  the development set, few-shot set or prompt template with a short change
  note. Every development run records the versions it used, so that any
  reported development result can be traced and re-run. New development
  families are checked automatically against the held-out manifest for
  near-duplicates. The check reports only that a conflict exists, not
  which held-out item matched; a flagged family is discarded or rewritten.
\item
  \textbf{Split by statement family.} A \emph{family} is a base
  statement and all its paraphrases, numeric variants and context
  variants. Families, not statements, are assigned to splits, to reduce cross-split paraphrase leakage. Family-ID and normalized exact-text checks alone cannot establish the absence of semantic near-duplicates; independent review is required.
\item
  \textbf{Context separation.} Where two or more institutional contexts
  exist, at least one context {[}TBD{]} appears only in the held-out set
  to test sensitivity to entity tables.
\item
  \textbf{No test tuning.} Prompts, few-shot examples, schema
  descriptions and decoding settings are chosen using
  development/validation data only. Few-shot exemplars are drawn from
  the development split. Before the confirmatory run, the complete
  configuration (model identifier, prompt template version, few-shot set
  version, schema version, decoding and retry settings) is \emph{frozen}
  and recorded; frozen versions are never edited, and a later change
  creates a new version that requires its own pre-specified run.
\item
  \textbf{One confirmatory pass.} The primary held-out evaluation is run
  once per frozen configuration (with the planned repetitions,
  Section\textasciitilde3.7). Held-out results are not used to revise
  prompts or examples for that study. If a defect is found afterwards,
  the fix and a re-run are reported as a deviation, not silently
  replacing the first result.
\item
  \textbf{Training-data contamination.} Statements are authored for this
  study and are not published before the evaluation. If the dataset is
  later shared, models released afterwards might have seen it; this is a
  stated limitation of any later re-use.
\item
  \textbf{Prompt-sensitivity.} Because SCHEDBench-type results suggest
  that wording and ordering can change model behaviour~\cite{sharma2026schedbench}
  (preprint, author-reported), the dataset includes
  paraphrase families, and consistency across paraphrases is reported
  (Section\textasciitilde3.9).
\end{itemize}

\hypertarget{llm-selection-and-equivalent-experimental-conditions}{%
\section{LLM selection and equivalent experimental
conditions}\label{llm-selection-and-equivalent-experimental-conditions}}

\textbf{Selection criteria (stated before any model is run).} (1)
Documented support for schema-constrained (structured) output; (2)
availability through an API under terms acceptable for the data used
(OQ-18); (3) a stable, versioned model identifier; (4) listed price and
rate limits compatible with the budget (OQ-13); (5) context size
sufficient for the entity context; (6) vendor diversity. Candidate
models are described in their vendor documentation~\cite{openai_gpt6_luna,google_gemini38_flash,anthropic_claude_docs}.
Open-weight/local models may be added if privacy or cost requires.

\textbf{Candidate selection remains open.} Earlier planning documents contain provisional model names; these are not treated here as verified availability, supported capabilities or approved purchases. Before experiments, official provider documentation must be checked and dated for exact model identifiers, structured-output support, prices, rate limits and data terms. Candidate selection must follow the stated criteria and an approved budget. No extraction benchmark has yet been run in this project. AI tools have assisted drafting and implementation; their use will be disclosed according to university policy, with the author responsible for checking claims and experimental records.

\textbf{Equivalent conditions.}

\begin{itemize}
\tightlist
\item
  Same constraint vocabulary and schema (with each
  provider\textquotesingle s schema-constraint mechanism;
  provider-specific schema subsets, if required, are documented and
  validated by the same backend validator).
\item
  Same statements, same institutional contexts, same minimal entity
  context format, same task instruction text (provider-specific wrapper
  changes only, logged).
\item
  Same number of repetitions per statement and model {[}TBD, e.g., 3--5;
  chosen from the pilot variance{]}.
\item
  Decoding settings as equivalent as each API allows (e.g.,
  lowest-variability setting available); where a setting cannot be
  matched (for example reasoning effort or sampling controls differ),
  the setting is recorded and the difference reported as a limitation.
  {[}TBD after checking each API.{]}
\item
  Same retry policy (bounded, e.g., at most one repair attempt for
  malformed output) and the same treatment of refusals and truncation
  (counted as failures).
\item
  Model identifiers and versions pinned; run dates recorded; runs for
  different models interleaved in time to reduce drift effects.
\item
  No fine-tuning (outside initial scope). Prompting factors (zero-shot
  versus few-shot, or presence of a vocabulary glossary) are decided on
  development data and fixed before the held-out run; they are not
  varied on test data.
\end{itemize}

\hypertarget{baselines}{%
\section{Baselines}\label{baselines}}

\begin{enumerate}
\def\labelenumi{\arabic{enumi}.}
\tightlist
\item
  \textbf{Manual-entry baseline (human, RQ3):} the structured forms of
  FR-010 produce the same constraint objects as the AI route; the same
  scheduling engine is used downstream.
\item
  \textbf{Rule-based extraction baseline (optional, RQ1):} a simple
  keyword/pattern parser for the vocabulary, to show what a non-LLM
  approach achieves on the same data. {[}TBD: include if time
  permits.{]}
\item
  \textbf{No-review versus post-review:} the raw LLM output is itself a
  reference point for the post-review result
  (Section\textasciitilde3.9.4).
\end{enumerate}

\hypertarget{metrics}{%
\section{Metrics}\label{metrics}}

All extraction measures will be reported overall and per language category with item counts. Small per-category estimates and intervals will be labelled exploratory. Entity-resolution evaluation includes Arabic and transliterated mentions. Scheduling checks include instructor eligibility, continuity across weekly occurrences, locks, and configured workload limits measured in elapsed teaching hours. Independent engineering checks are implemented; extraction and participant-study measures remain planned rather than measured research results.

\hypertarget{extraction-accuracy-primary-raw-output-held-out-set}{%
\subsection{Extraction accuracy (primary: raw output, held-out
set)}\label{extraction-accuracy-primary-raw-output-held-out-set}}

\begin{itemize}
\tightlist
\item
  \textbf{Exact semantic match (statement level):} proportion of
  statements where the predicted canonical constraint set equals the
  gold set (Section\textasciitilde3.3). \emph{Pre-specified primary
  metric.}
\item
  \textbf{Constraint-level precision, recall, F1:} a predicted
  constraint counts as true positive if semantically equal to a gold
  constraint (one-to-one matching); micro-averaged over the held-out
  set, with per-type breakdown.
\item
  \textbf{Field-level accuracies:} type; severity (hard/soft);
  entity-resolution accuracy (correct identifier among resolved
  mentions); parameter accuracy (days, intervals, thresholds, polarity);
  named-period resolution.
\item
  \textbf{Negation and polarity accuracy} on the negated subset.
\item
  \textbf{Paraphrase consistency:} for each family, the proportion of
  paraphrases yielding the same canonical set as the
  family\textquotesingle s base item.
\item
  \textbf{Run-to-run consistency:} across repetitions, the proportion of
  statements with identical outputs, and the spread of the primary
  metric.
\end{itemize}

\hypertarget{entity-handling}{%
\subsection{Entity handling}\label{entity-handling}}

\begin{itemize}
\tightlist
\item
  \textbf{Invented-entity rate:} raw outputs containing an identifier or
  mention treated as resolved that is not in the supplied context (per
  100 statements). Post-validation this should be zero by design
  (FR-016); the raw rate is still reported.
\item
  \textbf{Ambiguity handling:} on items with ambiguous entities, whether
  the output marks ambiguity with the correct candidates.
\end{itemize}

\hypertarget{clarification-and-unsupported-requests}{%
\subsection{Clarification and unsupported
requests}\label{clarification-and-unsupported-requests}}

\begin{itemize}
\tightlist
\item
  \textbf{Clarification trigger recall and precision:} on gold items
  needing clarification (missing value, ambiguous entity, undefined
  term, unclear severity), whether the system raised a clarification of
  the right kind; unnecessary clarifications on clear items.
\item
  \textbf{Question usefulness:} rated by annotators on a simple rubric
  (specific to the missing field; answerable; not leading) {[}TBD
  scale{]}.
\item
  \textbf{Unsupported handling:} proportion of unsupported items flagged
  as unsupported; \textbf{false activation} --- unsupported content that
  appears as an approved-looking constraint (should be 0 after the
  system\textquotesingle s gate, and is reported raw).
\item
  \textbf{Contradiction detection:} proportion of gold-contradictory
  items for which the contradiction is reported.
\item
  \textbf{Guessing rate:} items where the output supplies a value that
  the statement and context do not provide.
\end{itemize}

\hypertarget{raw-versus-post-review-rq1rq3}{%
\subsection{Raw versus post-review
(RQ1/RQ3)}\label{raw-versus-post-review-rq1rq3}}

\begin{itemize}
\tightlist
\item
  \textbf{Raw quality:} metrics above on the model output before human
  edits.
\item
  \textbf{Post-review quality:} the same metrics on the approved
  constraint set after review (in the human pilot and, for an upper
  bound, in a researcher-simulated review).
\item
  \textbf{Correction effort:} number of fields changed and constraints
  added/\allowbreak removed/\allowbreak rejected per statement (field-level diff between raw and
  approved); reviewer decision counts; review time.
\item
  \textbf{Residual error after review:} gold versus approved set, to
  show errors that review failed to catch.
\end{itemize}

\hypertarget{operational-metrics-rq4}{%
\subsection{Operational metrics (RQ4)}\label{operational-metrics-rq4}}

\begin{itemize}
\tightlist
\item
  \textbf{Latency} per extraction (request to validated proposal),
  reported as median and upper percentiles, noting network and provider
  load effects.
\item
  \textbf{Tokens and cost:} input/output tokens and cost computed from
  the pricing in force on the run date, stored with the source and date.
  \textbf{Cost per successful extraction} = total cost divided by the
  number of statements with an exact semantic match (raw), reported
  alongside cost per statement.
\item
  \textbf{Reliability:} rates of refusals, malformed or truncated
  output, rate-limit errors, retries, and failures.
\end{itemize}

\hypertarget{395a-secondary-downstream-equivalence-on-small-fixtures}{%
\subsection{Secondary: downstream equivalence on small
fixtures}\label{395a-secondary-downstream-equivalence-on-small-fixtures}}

For a small subset of items with fixture term data, the extracted
constraint set and the gold set can each be compiled into the solver
model and compared by whether they accept the same set of candidate
timetables (a set of fixed timetables, each checked by fixing every
assignment variable, as in Section\textasciitilde3.11). This is reported
separately as a secondary measure in the spirit of solution-based
judging used in constraint-modelling benchmarks~\cite{michailidis2025cpbench}. It can only show equivalence on the tested timetables, not in
general, and is never used in place of the field-level scoring above.

\hypertarget{scheduling-and-validation-metrics}{%
\subsection{Scheduling and validation
metrics}\label{scheduling-and-validation-metrics}}

Defined in Section\textasciitilde3.11.

\hypertarget{human-in-the-loop-comparison-rq3--proposed-pilot}{%
\section{Human-in-the-loop comparison (RQ3) --- proposed
pilot}\label{human-in-the-loop-comparison-rq3--proposed-pilot}}

This study is conditional on participant access, ethics requirements and feasibility. If conducted, participants may use their preferred supported language category; the category will be recorded per task. No participant count or approval is assumed.

\textbf{Design.} Within-subject, counterbalanced crossover: each
participant completes one task list with manual forms and an equivalent
task list with AI-assisted entry and review; order and task-list
assignment are counterbalanced. Task lists contain the same number of
requirements and the same mix of categories. They are drawn from a
separate task pool that is disjoint from the held-out set and not used
for model development, so the pilot does not expose held-out items, and
they are matched by pilot difficulty ratings.
{[}TBD: number of requirements per list.{]}

\textbf{Participants.} Preferably people with timetabling or
administrative scheduling experience (staff, teaching assistants),
otherwise students as a stated proxy. \textbf{No participant number is
assumed.} The number will be chosen with the supervisor on feasibility
grounds and reported as a pilot if small; a power calculation is not
meaningful without effect-size information, so none is claimed. {[}TBD:
n; recruitment channel.{]}

\textbf{Procedure.} Briefing and consent; short training on both
interfaces; timed tasks; post-task questionnaires; brief interview. The
experiment mode (FR-050) records start/end times, edits and corrections;
no behavioural data is captured without a recorded consent. Instruments
for perceived workload and usability are standard questionnaires to be
selected {[}TBD, e.g., a workload scale and a usability scale{]}.

\textbf{Educational and administrative task context.} Tasks will include protecting access to a required lecture and laboratory, entering staff availability, requesting suitable laboratory facilities, and expressing a preference about gaps or late teaching. Cases will use the supported group-based course scope; exam-spacing and individual degree-completion tasks are excluded unless their extensions are approved. After each task, the participant will explain what the approved rule permits and forbids. Understanding will be scored against a predefined rubric, and interviews will ask whether the displayed interpretation and conflict feedback are actionable. This complements speed and accuracy measures; it does not measure student learning outcomes.

\textbf{Outcomes.} Task completion time; number of corrections (AI
condition) and form-edit counts (manual condition); accuracy of the
final approved constraint set against gold in both conditions (so that
speed is not credited if correctness falls); self-reported workload and
preference; correctness of the explanation of the approved rule; qualitative comments about clarity and actionable feedback.

\textbf{Analysis.} Paired comparison per participant (median
differences, Wilcoxon signed-rank or equivalent) with effect sizes and
confidence intervals, reported descriptively when n is small; no claim
of broad usability improvement from a pilot. Order and learning effects
are examined by comparing orders.

\textbf{Fairness.} Both conditions use the same scheduling engine,
vocabulary and downstream pipeline, so differences arise from how
constraints are entered.

\hypertarget{scheduling-validation-and-system-behaviour-evaluation}{%
\section{Scheduling, validation and system-behaviour
evaluation}\label{scheduling-validation-and-system-behaviour-evaluation}}

\textbf{Instances.} (a) Small hand-built instances, each designed to
exercise one constraint type; (b) a synthetic department generator with
documented parameters (activities, rooms, instructors, groups, grid) and
recorded seeds; (c) optionally, a small benchmark-derived instance with
\emph{explicitly mapped semantics}, noting that the ITC 2019 model
contains features (weeks, student sectioning, many distribution types)
outside the core scope~\cite{muller2025itc}. {[}TBD: scale.{]} Tested
scale is reported with every result.

\textbf{Measures.} Assignment completeness; hard violations found by the
independent validator (expected zero for published schedules); soft
penalties computed by the validator (and compared with the solver
objective as an integrity check); solver outcome distribution
(optimal/feasible/infeasible/unknown/model-invalid); runtime and time
limits; objective definition; seeds and worker count. Runs are repeated
across seeds {[}TBD count{]} because multi-worker search need not be
deterministic (assumption A-09); variation is reported.

\textbf{Not attributing solver differences to AI.} When the manual and
AI-assisted conditions yield different approved constraint sets,
schedule differences are reported as \emph{consequences of different
inputs}, and totals are compared only under a common objective
definition (FR-036).

\textbf{Verifying the validator.} (1) \emph{Seeded-violation (mutation)
tests:} for each supported rule, valid timetables are mutated to
introduce exactly one violation of that rule; the validator must report
each, with the correct affected activities (target 100 \% by
construction; says nothing about unseeded cases). (2) \emph{Differential
testing:} randomly perturbed timetables are checked by both the
validator and the solver model. For the solver model, every assignment
variable is fixed to the timetable\textquotesingle s values by an
equality constraint (not a solver hint, which only guides the search and
could let the solver repair an invalid timetable) and the model is
solved for feasibility: a feasible result means the solver model accepts
the timetable, INFEASIBLE means it rejects it, and UNKNOWN is recorded
as inconclusive and is never counted as agreement. Disagreements between
the two are investigated. This reduces, but does not eliminate, the risk
that the same misconception is built into both
(Section\textasciitilde3.15). (3) \emph{Coverage tests:} inputs with
missing group or instructor data must produce coverage statements, not
blanket validity.

\textbf{System-behaviour tests (acceptance).} Scripted tests of SRS
acceptance criteria AC-01--AC-21, for example: an unapproved proposal
never appears in a snapshot; an unsupported request creates no
constraint; a tiny time limit on a feasible instance yields UNKNOWN or
FEASIBLE and never INFEASIBLE; an edit that creates a violation blocks
publication; stale scenarios are flagged; exports carry scenario and
input version. The additional criteria AC-16--AC-21 will be tested through paired concurrent saves and an input change during publication; interruption and reconnect cases; duplicate requests and exhausted provider retries; partial results and completion/cancellation races; validator failures and missing coverage; and expired sessions or revoked roles. Each case will specify the expected state, preserved records, user-visible explanation and permitted recovery action. Evidence will include whether publication remained blocked when required and whether restricted data stayed inaccessible. Pass/fail and any unresolved failures will be reported against each criterion; these checks do not establish production reliability or improved learning outcomes. No statistics are used.

\hypertarget{analysis-methods-and-uncertainty}{%
\section{Analysis methods and
uncertainty}\label{analysis-methods-and-uncertainty}}

\begin{itemize}
\tightlist
\item
  \textbf{Pre-specification.} The primary metric (raw exact semantic
  match on the held-out set), the model comparison procedure and the
  exclusion rules are recorded before the held-out run.
\item
  \textbf{Intervals.} Proportions are reported with Wilson 95 \%
  intervals; clustered bootstrap over statement families for
  micro-averaged scores and F1.
\item
  \textbf{Model comparison.} Models are compared on the same statements:
  paired comparison of per-statement correctness using
  McNemar\textquotesingle s exact test (or a family-clustered
  alternative), with effect sizes and Holm correction across the number
  of pairwise comparisons; with three models there are three pairs.
  Statistical significance is secondary to interval width and practical
  difference.
\item
  \textbf{Repetitions.} Each model is run {[}TBD{]} times per statement.
  Reported: mean and range of the primary metric across runs,
  per-statement stability; the primary result is the mean across runs,
  and a majority-vote result is reported separately only if a deployment
  would use it.
\item
  \textbf{Latency and cost.} Descriptive (median, percentiles), with the
  caveat that latency depends on provider load.
\item
  \textbf{Category analyses} are exploratory unless sample sizes justify
  them.
\item
  \textbf{Reporting deviations.} Any deviation from the plan is listed
  with a reason.
\item
  \textbf{Model choice rule.} Following the project specification, the
  least expensive model that meets agreed accuracy and clarification
  targets on held-out data is preferred. The targets are \textbf{not yet
  set} (OQ-10) and will be proposed after the pilot and agreed before
  the held-out run.
\end{itemize}

\hypertarget{reproducibility-records}{%
\section{Reproducibility records}\label{reproducibility-records}}

The complete research protocol requires the following record for every extraction, solver run and experiment; this is a target specification rather than a claim that provider instrumentation already exists. Required fields are
(FR-023, FR-024, FR-025; NFR-016): dataset and split version (with the
development-set change log and the held-out manifest hash); annotation
guide version; few-shot set version; institutional context version; prompt template and schema
versions; provider and model identifier; request date/time; decoding
settings; context sent; raw response; retries and errors; token counts;
latency; cost with its pricing source and date; input snapshot version
and hash; solver version, parameters, seed and time limit; runtime and
outcome; validator version. Experiment configuration files are
version-controlled; analysis scripts are archived with the results.
Because provider models may change or be withdrawn, raw outputs are
archived so that scoring can be re-run even if the model cannot.

\hypertarget{ethics-privacy-and-consent}{%
\section{Ethics, privacy and consent}\label{ethics-privacy-and-consent}}

\begin{itemize}
\tightlist
\item
  \textbf{Data minimisation.} Only statement text and the minimal entity
  context (names, aliases, identifiers, named periods) are sent to an
  LLM provider; no student-level personal data is sent (FR-014,
  NFR-006). Synthetic or anonymised data are used wherever possible
  (A-10).
\item
  \textbf{Provider terms.} The data-use and retention terms of each API
  provider are reviewed before use {[}TBD, OQ-18{]}; settings that
  exclude data from model training are used where offered.
\item
  \textbf{Human participants (if the pilot proceeds).} Institutional
  ethics or research-approval requirements are checked first {[}TBD,
  OQ-12{]}; participation is voluntary with informed consent covering
  purpose, data collected, retention, and the right to withdraw; data
  are pseudonymised; the pilot involves no sensitive personal data;
  results are reported in aggregate.
\item
  \textbf{Real university data.} Used only with written permission from
  the data owner and in compliance with applicable data-protection law
  and university policy {[}TBD, OQ-16, OQ-17{]}; otherwise not used.
\item
  \textbf{Dataset sharing.} Statements derived from real staff
  communications are shared only if anonymised and permitted.
\item
  \textbf{Honest reporting.} Negative or null findings are reported;
  claims are limited to the sample and settings.
\end{itemize}

\hypertarget{threats-to-validity}{%
\section{Threats to validity}\label{threats-to-validity}}

\textbf{Construct validity.} Gold annotations encode one interpretation;
ambiguous statements may have several reasonable readings (mitigated by
the guide, a second annotator on a subset, and by treating "needs
clarification" as a gold outcome). Exact semantic match is strict;
looser metrics are reported alongside. \textbf{Internal validity.}
Prompt tuning on development data can overfit; leakage between splits
via paraphrases (mitigated by family-level splits and frozen prompts);
provider updates changing model behaviour during the study (pinned
identifiers, interleaved runs, logged dates); non-equivalent decoding
controls across providers; researcher bias because the author builds the
system, writes the statements and annotates them (second annotator,
pre-specification, held-out discipline). \textbf{External validity.}
Synthetic statements and fictional departments may be easier or
different from real administrators\textquotesingle{} language;
Four language categories with potentially unequal sample sizes; Egyptian Arabic dialect variation not systematically covered; author-written Arabic may not represent administrative usage; one department and one term; closed vocabulary; results
may not transfer to other institutions, languages or models released
later; pilot participants may not represent schedulers.
\textbf{Conclusion validity.} Small sample sizes, repeated testing of
several models, and dependence among paraphrases (handled by clustered
analysis and conservative reporting). \textbf{Validator independence.}
One person implementing both solver encoding and validator can build the
same misunderstanding into both; mitigations are the written semantics,
seeded tests, differential tests and, if feasible, review by another
person. Independence is relative to the solver encoding, not absolute.
\textbf{Conflict of interest.} The drafting assistant is a Claude model
and Claude is a candidate; mitigated by pre-specified selection rules
and by reporting all candidates\textquotesingle{} results.

\hypertarget{feasibility-risks-and-limitations}{%
\section{Feasibility risks and
limitations}\label{feasibility-risks-and-limitations}}

\begin{longtable}[]{@{}p{\dimexpr(\linewidth-6\tabcolsep)/3\relax}p{\dimexpr(\linewidth-6\tabcolsep)/3\relax}p{\dimexpr(\linewidth-6\tabcolsep)/3\relax}@{}}
\toprule\noalign{}
Risk & Effect & Mitigation \\
\midrule\noalign{}
\endhead
\bottomrule\noalign{}
\endlastfoot
Annotation effort exceeds capacity & Smaller dataset, wider intervals &
Choose Option A; prioritise held-out quality; report intervals \\
No second annotator & No agreement estimate & State limitation; use
guide-based self-consistency re-annotation after a delay as a weak
substitute, labelled as such \\
Participants unavailable & No RQ3 evidence & Report RQ3 as not answered
or as a very small pilot \\
API cost or model withdrawal & Fewer models or runs & Budget (OQ-13);
archive outputs; adapter allows substitution \\
Solver scale unknown & Limited scheduling claims & Early spike; report
tested scale \\
Multilingual annotation and instructor assignment & Delay & Agree category coverage and test instructor variables early; extensions
only after core evaluation \\
Real data unavailable & Less realistic statements & Synthetic data;
disclose \\
\end{longtable}

\textbf{Limitations stated in advance.} Results will concern a closed
vocabulary, English/formal Arabic/Egyptian Arabic/mixed-language statements, synthetic one-department-sized instances, institution-specific workload values and
specific model versions on specific dates. Validity of a schedule is
relative to the supported rules and supplied data. Extraction accuracy
on the benchmark does not imply correct schedules in deployment, and
valid schedules do not imply correct extraction.
